The proof of the full-vector exact-count theorem is coordinatewise.  Restrict it to \(k\in A\).  The assumption \(C_{C,k}/C\to\alpha_k>0\) on the active face alone ensures that every active assignment-cell count diverges, that the expansion (8) holds on \(A\), and that the triangular-array maximum is \(O(C^{-1/2})\).  The active conditional covariance converges to
\[
 \operatorname{diag}(\tau_k^2/\alpha_k:k\in A)=\Sigma_{CR,A}.
\]
The nondegeneracy assumption makes this matrix positive definite and validates active-coordinate studentization.  Applying the same characteristic-function and empirical-second-moment arguments on the finite set \(A\) proves both displayed convergences.  No step refers to a cell with \(k\notin A\), so a zero inactive share creates no asserted estimator or limit.